\documentclass[12pt]{article}
\usepackage{amsmath, amssymb}
\usepackage{geometry}
\usepackage{booktabs}
\usepackage{parskip}

\geometry{margin=1in}

\title{HW3 --- Question 3c\\{\large ORF 309}}
\author{}
\date{}

\begin{document}
\maketitle

%---------------------------------------------------------------
\section*{Problem Statement}

A bag contains coins: $\frac{2}{3}$ are fair (P(H) = 0.5) and $\frac{1}{3}$ are Canadian
(P(H) = 0.65). One coin is drawn uniformly at random and flipped twice independently.
Given that exactly 1 tail was observed, what is the probability the coin is Canadian?

%---------------------------------------------------------------
\section*{Setup and Notation}

Define the following events:
\begin{align*}
F &= \{\text{coin drawn is fair}\}, \quad &&P(F) = \tfrac{2}{3} \\
C &= \{\text{coin drawn is Canadian}\}, \quad &&P(C) = \tfrac{1}{3} \\
E &= \{\text{exactly 1 tail in 2 flips}\}
\end{align*}

We want to compute $P(C \mid E)$.

%---------------------------------------------------------------
\section*{Step 1: Compute $P(E \mid F)$ and $P(E \mid C)$}

Given the coin type, the two flips are independent and identically distributed.
So the number of tails $T \mid \text{coin type} \sim \mathrm{Binomial}(2,\, p_T)$,
where $p_T$ is the probability of tails on a single flip.

\subsection*{Fair coin: $p_T = 0.5$}

The two sequences giving exactly 1 tail are HT and TH:

\begin{align*}
P(E \mid F)
&= P(\text{HT} \mid F) + P(\text{TH} \mid F) \\
&= (0.5)(0.5) + (0.5)(0.5) \\
&= 0.25 + 0.25 \\
&= \boxed{0.5}
\end{align*}

Equivalently, using the Binomial formula with $n=2$, $k=1$, $p_T = 0.5$:
\[
P(E \mid F) = \binom{2}{1}(0.5)^1(0.5)^1 = 2 \cdot 0.25 = 0.5
\]

\subsection*{Canadian coin: $p_T = 0.35$}

\begin{align*}
P(E \mid C)
&= P(\text{HT} \mid C) + P(\text{TH} \mid C) \\
&= (0.65)(0.35) + (0.35)(0.65) \\
&= 0.2275 + 0.2275 \\
&= \boxed{0.455}
\end{align*}

Equivalently, using the Binomial formula with $n=2$, $k=1$, $p_T = 0.35$:
\[
P(E \mid C) = \binom{2}{1}(0.35)^1(0.65)^1 = 2 \cdot (0.35)(0.65) = 2 \cdot 0.2275 = 0.455
\]

%---------------------------------------------------------------
\section*{Step 2: Compute $P(E)$ via the Law of Total Probability}

Since $F$ and $C$ are mutually exclusive and exhaustive (every coin is either fair
or Canadian):
\[
P(E) = P(E \mid F)\,P(F) + P(E \mid C)\,P(C)
\]

Substituting:
\begin{align*}
P(E) &= (0.5)\!\left(\frac{2}{3}\right) + (0.455)\!\left(\frac{1}{3}\right) \\[6pt]
     &= \frac{1}{3} + \frac{0.455}{3} \\[6pt]
     &= \frac{1 + 0.455}{3} \\[6pt]
     &= \frac{1.455}{3} \\[6pt]
     &\approx 0.485
\end{align*}

%---------------------------------------------------------------
\section*{Step 3: Apply Bayes' Formula}

\[
P(C \mid E) = \frac{P(E \mid C)\,P(C)}{P(E)}
\]

Substituting:
\begin{align*}
P(C \mid E)
&= \frac{(0.455)\!\left(\dfrac{1}{3}\right)}{\dfrac{1.455}{3}} \\[10pt]
&= \frac{\dfrac{0.455}{3}}{\dfrac{1.455}{3}} \\[10pt]
&= \frac{0.455}{1.455} \\[6pt]
&= \frac{91}{291} \\[6pt]
&\approx \boxed{0.313}
\end{align*}

\textit{Exact fraction derivation:}
\[
0.455 = \frac{91}{200}, \qquad 1.455 = \frac{291}{200}
\qquad \Longrightarrow \qquad
\frac{0.455}{1.455} = \frac{91/200}{291/200} = \frac{91}{291}
\]

%---------------------------------------------------------------
\section*{Summary}

\begin{center}
\begin{tabular}{lcc}
\toprule
 & Fair ($F$) & Canadian ($C$) \\
\midrule
Prior $P(\cdot)$ & $2/3 \approx 0.667$ & $1/3 \approx 0.333$ \\
$P(E \mid \cdot)$ & $0.5$ & $0.455$ \\
$P(E \mid \cdot)\cdot P(\cdot)$ & $0.333$ & $0.152$ \\
\midrule
Posterior $P(\cdot \mid E)$ & $0.687$ & $\mathbf{0.313}$ \\
\bottomrule
\end{tabular}
\end{center}

\medskip
\textbf{Interpretation.} The prior probability of a Canadian coin was $33.3\%$.
After observing exactly 1 tail, the posterior drops slightly to $31.3\%$.
This makes sense: exactly 1 tail in 2 flips is slightly \emph{more} consistent
with a fair coin ($P(E \mid F) = 0.5$) than with a Canadian coin
($P(E \mid C) = 0.455$), so the evidence nudges us weakly toward fair --- but
not by much, since the likelihoods are close.

\end{document}
